% =============================================================================
% Glossary
% =============================================================================

\section*{Glossary}
\addcontentsline{toc}{section}{Glossary}

\begin{description}[style=nextline,leftmargin=0pt]

\item[VarFrame]
The primary N-dimensional data structure in \itc{}. Stores named dimensions
(coordinate arrays) and named variables (NumPy arrays). Immutable by default.
Carries Provenance and History metadata.

\item[UncFrame]
Mirror of a VarFrame storing standard uncertainties for each variable. Created
automatically when uncertainties are assigned. Indexed identically to its
parent VarFrame. Structurally a VarFrame of standard uncertainties.

\item[HistoryFrame]
Mirror of a VarFrame encoding the origin tag of every value in every variable.
Each entry takes one of three integer values: \code{0} for values loaded
directly from the source files, \code{+1} for values that were interpolated or
computed, and \code{-1} for values that were extrapolated. Enables visual and
programmatic inspection of which parts of a dataset are measured ground truth
versus model-derived. Reductions follow a worst-case rule: the result tag is
\code{-1} if any contributor is \code{-1}, else \code{+1} if any is \code{+1},
else \code{0}.

\item[Provenance]
The static, immutable record of where the data came from: source files,
content hashes, user, creation timestamp, \itc{} version, and operating mode.
Travels with the VarFrame and is preserved in every export format.

\item[History]
The cumulative, ordered, append-only record of every operation applied to a
VarFrame since its creation. Each entry has a unique sequential index, the
operation name and arguments, an optional user comment, a timestamp, and a
state hash. Distinct from Provenance (origin) by purpose: History describes
what happened to the data.

\item[Pipeline]
A serializable subrange of a History, exportable as a \code{.itc\_pipe} file
and reusable on any compatible VarFrame.

\item[Axis]
A coordinate-frame object stored in \code{core/axes.py}, carrying a name and a
$3\times 3$ rotation matrix relative to the canonical AIAA body axis. Custom
axes (stability, wind, tunnel, propeller-disk) are first-class objects;
\code{db.rotate(target\_axis)} applies the rotation to detected vector groups
in the VarFrame and propagates uncertainty using the rotation matrix as the
Jacobian.

\item[compute]
The \itc{} operation that derives new variables from string equations, with
automatic uncertainty propagation when applicable, including correlation
terms when a covariance matrix has been declared.

\item[interpolate]
The \itc{} operation that resamples variables onto new coordinate values along
one or more dimensions, optionally translating the sweep axis (e.g.\
$\alpha \to C_L$). By default, coordinate values that already exist in the
target dimension are preserved and not recomputed; pass \code{override=True}
to force recomputation.

\item[fitmodel / fitvalue]
Operation pair that replaces a dimension by polynomial coefficient labels of
the form $\alpha^0, \alpha^1, \dots, \alpha^N$ (\code{fitmodel}) and evaluates
the resulting polynomial fit at user-specified coordinates (\code{fitvalue}).

\item[squeeze / expand / concat]
Dimension-shape operations: \code{squeeze} removes unit-cardinality
dimensions; \code{expand} adds a new dimension by broadcasting;
\code{itc.concat} concatenates compatible VarFrames along a shared dimension
with non-overlapping coordinate values.

\item[combine]
Explicit arithmetic combination of two VarFrames. Replaces overloaded Python
arithmetic operators to keep every numerical combination traceable. Example:
\code{db.combine(other, op="mean")} returns a new VarFrame whose variables
are the elementwise mean of the inputs, with uncertainty propagated under the
declared correlation structure.

\item[pproc]
Processor object in \itc{}. Receives a VarFrame and returns an enriched
VarFrame. Encapsulates a complete analysis workflow including corrections,
derivations, and uncertainty propagation. Conforms to the \code{Processor}
typing protocol.

\item[.itceq]
\itc{} equation file. A TOML-structured text file that defines metadata,
constants, uncertainties, equations, and corrections for a campaign. Reusable
across tunnels and rigs. Version-controllable.

\item[.itc]
\itc{} native binary format. A ZIP archive containing variable arrays
(\code{.npz}), dimension definitions, uncertainty arrays, history-tag arrays,
provenance, history, and metadata as JSON. Designed for long-term archival.

\item[.itc\_pipe]
\itc{} pipeline file. Human-readable serialization of a sequence of operations
extracted from a VarFrame's History. Reusable across compatible VarFrame
objects.

\item[.itc\_surr]
\itc{} surrogate file. Binary serialization of a fitted surrogate model with
metadata. Optionally embeds the source VarFrame for full self-contained audit
(off by default).

\item[manifest]
A tabular export that maps each source file to its dimension coordinate
values. Generated after loading and before any data manipulation. Acts as the
primary audit checkpoint between data acquisition and analysis.

\item[diagnostics]
A structured inspection report for a VarFrame that identifies missing
combinations, non-finite values, partial variables, and per-variable
statistical summaries. Prints to terminal and returns an inspectable
\code{DiagnosticsReport} object.

\item[DataVis]
The data visualization object in \itc{}, created via \code{itc.datavis()}.
Holds figure-level style settings and exposes single-plot and multi-plot
generation backed by matplotlib (default) or plotly (interactive option).

\item[ItcFigure]
Figure wrapper returned by DataVis, exposing \code{export()},
\code{add\_legend()}, \code{annotate()}, and a \code{.mpl} escape hatch to
the underlying \code{matplotlib.Figure}.

\item[datapoint mode]
An operating state in which a VarFrame has only one synthetic dimension
named \code{datapoint} indexing rows in acquisition order. Reached by loading
without specifying \code{dims}. Always recoverable into a structured VarFrame
via \code{db.pivot()}.

\item[operating mode]
Either \code{production} (default; Provenance and History are mandatory and
exports are unrestricted) or \code{draft} (History is recorded only when the
user requests it explicitly; export requires \code{allow\_draft=True}). Mixing
modes in binary operations is forbidden: \code{db.promote()} or
\code{db.demote()} must be called explicitly first.

\item[EUB]
Engineering Uncertainty Band. The uncertainty band on an engineering decision
variable (e.g.\ trim deflection, takeoff distance, static margin) propagated
from input uncertainties through an \code{itc.aerospace} computation. Native
output of \itc{}'s uncertainty machinery, not a separate analysis mode.

\end{description}

\clearpage
